\subsection{生成多项式中的 Fronsdal 恒等式}\label{sec:ghs-fronsdal-higher-spin}

分量方程已经给出自由高自旋场的结构。辅助多项式将所有自旋的组合系数同时固定下来，并允许逐行核查规范变化与 Bianchi 恒等式。

此前把高自旋理解为张量丛的纤维表示，也讨论了共形 Killing 张量产生的对称算子。然而在量子场论里，\emph{自旋为 $s$ 的局域场}和\emph{传播自旋为 $s$ 的无质量粒子}不是同一句话。后者还必须说明哪些分量是约束、哪些由规范变换识别，以及留下的自由度如何承载质量为零的 Poincar\'e 表示。在平直洛伦兹时空，这些问题能够对任意整数 $s$ 作完整的微分算子计算。

本节使用 $x^0=ct$ 进行量纲核对；令 $c=1$ 后回到前面时空章节的自然单位记法。令时空维数为 $D\ge4$，度规取 $\eta_{\mu\nu}=\operatorname{diag}(-1,1,\ldots,1)$，笛卡尔坐标 $x^\mu$，$x^0=ct$。令 $\varphi_{\mu_1\cdots\mu_s}(x)$ 为实对称秩 $s$ 的协变张量。$\varphi$ 在平直空间没有预设无迹条件；对 $s\ge4$ 只要求\emph{双迹}消失。它与此前作为复合算符的对称无迹高自旋守恒流 $\mathcal J^{(s)}$ 不同。

为避免每一行重新排列数十个指标，先定义一个纯代数的辅助变量 $u^\mu$。取 $u^\mu$ 为无量纲记账变量，把张量编码为
\begin{equation}\label{eq:ghs-r7-fronsdal-field-polynomial}
 \Phi(x,u)=\frac1{s!}\varphi_{\mu_1\cdots\mu_s}(x)
 u^{\mu_1}\cdots u^{\mu_s}.
\end{equation}
任意关于 $u$ 的次数为 $s$ 的齐次多项式与对称张量一一对应；$u$ 不是新的时空坐标，也不是物理动量。定义五个微分算子
\begin{equation}\label{eq:ghs-r7-fronsdal-operators}
 U=u^\mu\partial_\mu,\quad
 D_u=\partial_u\!\cdot\!\partial_x,\quad
 T_u=\partial_u\!\cdot\!\partial_u,\quad
 \square=\eta^{\mu\nu}\partial_\mu\partial_\nu,\quad
 u^2=\eta_{\mu\nu}u^\mu u^\nu.
\end{equation}
其中 $U$ 将张量的阶数增加一，$D_u$ 将阶数减少一，$T_u$ 实施一次迹缩并；而 $\square$ 只作用于 $x$。按照上述有 $1/s!$ 的约定，$T_u\Phi$ 恰编码通常的单迹张量，不会额外出现不明来源的 $s(s-1)$。

\begin{lemma}[高自旋微分算子的交换代数]\label{lem:ghs-r7-fronsdal-algebra}
对任意 $x,u$ 上的多项式系数光滑函数，成立
\begin{equation}\label{eq:ghs-r7-fronsdal-commutators}
 [D_u,U]=\square,\qquad [T_u,U]=2D_u,\qquad
 [D_u,u^2]=2U,\qquad [T_u,u^2]=4N_u+2D,
\end{equation}
其中 $N_u=u\cdot\partial_u$ 是 $u$ 的次数算子。并且 $[D_u,T_u]=0$、$[U,\square]=0$。
\end{lemma}
\begin{proof}
第一个等式只需让 $D_uU-UD_u$ 作用于测试多项式 $F(x,u)$：
\begin{align*}
 D_u(UF)&=\partial_{u^\mu}\partial^\mu_x
 (u^\nu\partial_\nu F)
 =\partial^\mu\partial_\mu F
 +u^\nu\partial_\nu D_uF,\\
 U(D_uF)&=u^\nu\partial_\nu D_uF.
\end{align*}
两式相减即为 $\square F$。对第二式，先计算
$\partial_{u^\mu}(u^\nu\partial_\nu F)=\partial_\mu F+u^\nu\partial_\nu\partial_{u^\mu}F$，再次对 $u$ 求导得到 $2D_uF$。第三式来自 $\partial_{u^\mu}u^2=2u_\mu$。第四式由 $T_u(u^2F)$ 展开：两次导数均打在 $u^2$ 上时得到 $2DF$，一阶导数分别打在 $u^2$ 和 $F$ 上时得到 $4N_uF$。全部关系无需指定 $s$。
\end{proof}

把无质量自旋 $s$ 的二阶张量方程压缩为以下\emph{弗朗斯达尔算子}：
\begin{equation}\label{eq:ghs-r7-fronsdal-operator}
 \mathcal F[\Phi]=\square\Phi-U D_u\Phi
 +\frac12U^2T_u\Phi.
\end{equation}
为了确定其意义，不能仅凭形式像波动方程就宣称它是正确的动力学算子。令规范参数 $\varepsilon(x,u)$ 为次数 $s-1$ 的齐次多项式，定义
\begin{equation}\label{eq:ghs-r7-fronsdal-gauge}
 \delta_\varepsilon\Phi=U\varepsilon,\qquad
 T_u\varepsilon=0.
\end{equation}
在张量指标语言中这是 $\delta\varphi_{\mu_1\cdots\mu_s}=s\partial_{(\mu_1}\varepsilon_{\mu_2\cdots\mu_s)}$；圆括号表示除以排列数的完全对称化。

\begin{theorem}[任意整数自旋的平直规范恒等式]\label{thm:ghs-r7-fronsdal-gauge}
对任意次数 $s-1$ 的 $\varepsilon$，有
\begin{equation}\label{eq:ghs-r7-fronsdal-gauge-variation}
 \mathcal F[U\varepsilon]=\frac12U^3T_u\varepsilon.
\end{equation}
因此无迹规范参数条件 $T_u\varepsilon=0$ 使 $\mathcal F[\Phi]=0$ 对规范变换 \cref{eq:ghs-r7-fronsdal-gauge} 不变。
\end{theorem}
\begin{proof}
先展开三个项并将 $D_u$ 或 $T_u$ 向右移：
\begin{align*}
 \mathcal F[U\varepsilon]
 &=\square U\varepsilon-U D_uU\varepsilon
   +\frac12U^2T_uU\varepsilon\\
 &=U\square\varepsilon
   -U(UD_u+\square)\varepsilon
   +\frac12U^2(UT_u+2D_u)\varepsilon\\
 &=\bigl(U\square-U\square\bigr)\varepsilon
   +\bigl(-U^2D_u+U^2D_u\bigr)\varepsilon
   +\frac12U^3T_u\varepsilon.
\end{align*}
前两组抵消后只剩参数的迹。另外 $T_u^2(U\varepsilon)=2D_uT_u\varepsilon+UT_u^2\varepsilon=0$，所以规范变换也保持场的双迹约束。这里的 $1/2$ 与三次对称梯度来自统一的生成函数约定，任何 $s$ 均可使用。
\end{proof}

自旋一时 $T_u\varepsilon=0$ 是自动的，$\mathcal F$ 恢复 Maxwell 势方程；自旋二时参数为向量，迹同样自动为零，而 $\mathcal F$ 与线性化 Ricci 张量的组合有关。自旋 $s\ge3$ 时无迹条件才真正限制规范参数，这不是为了简化推导而人为选择的规范。

方程对规范变换不变，还不足以保证能由局域二次作用量导出。为了让规范方向在变分中成为零方向，需要构造与 $\mathcal F$ 相应的 Einstein 型张量，并证明它的散度满足适当恒等式。

\begin{proposition}[弗朗斯达尔 Bianchi 恒等式]\label{prop:ghs-r7-fronsdal-bianchi}
对任意多项式 $\Phi$，令 $\mathcal F$ 如 \cref{eq:ghs-r7-fronsdal-operator}，则
\begin{equation}\label{eq:ghs-r7-fronsdal-bianchi}
 \left(D_u-\frac12UT_u\right)\mathcal F[\Phi]
 =-\frac14U^3T_u^2\Phi.
\end{equation}
当 $s\le3$ 时 $T_u^2\Phi$ 自动为零；当 $s\ge4$ 时，要求 $T_u^2\Phi=0$ 可以消除右端。
\end{proposition}
\begin{proof}
依次使用 \cref{eq:ghs-r7-fronsdal-commutators}，并注意 $[D_u,T_u]=0$，得到
\begin{align*}
 D_u\mathcal F
 &=\square D_u\Phi-(UD_u+\square)D_u\Phi
  +\frac12(U^2D_u+2U\square)T_u\Phi\\
 &=-UD_u^2\Phi+\frac12U^2D_uT_u\Phi+U\square T_u\Phi.
\end{align*}
对后一个项进行迹缩并，需要先推导 $[T_u,U^2]=4UD_u+2\square$。于是
\begin{align*}
 T_u\mathcal F
 &=\square T_u\Phi-(UT_u+2D_u)D_u\Phi
 +\frac12(U^2T_u+4UD_u+2\square)T_u\Phi\\
 &=2\square T_u\Phi+UD_uT_u\Phi-2D_u^2\Phi
 +\frac12U^2T_u^2\Phi.
\end{align*}
把第二个式子乘 $U/2$ 后从第一个式子中减去，含 $D_u^2\Phi$、$D_uT_u\Phi$、$\square T_u\Phi$ 的项分别抵消，留下 $-U^3T_u^2\Phi/4$。
\end{proof}

定义 Einstein 型算子 $\mathcal G[\Phi]=\mathcal F[\Phi]-\tfrac14u^2T_u\mathcal F[\Phi]$。对称张量的配对取
\[
 \langle\Psi,\Phi\rangle_s
 =\frac1{s!}\int_{\RR^{1,D-1}}
 \overline{\varphi_{\mu_1\cdots\mu_s}}
 \psi^{\mu_1\cdots\mu_s}\,\dd^Dx,
\]
它是洛伦兹签名下的不定 Hermitian 配对，在实场子空间上成为对称双线性型；它不是正定 Hilbert 内积。形式变分时先假设场和变分紧支撑，边界积分因此为零。取保持 SI 量纲的常数 $\mathcal C_s$，定义
\begin{equation}\label{eq:ghs-r7-fronsdal-action}
 I_s[\Phi]=\frac{\mathcal C_s}{2}\langle\mathcal G[\Phi],\Phi\rangle_s,
 \qquad T_u^2\Phi=0.
\end{equation}
由于 $x^0=ct$ 及 $[\square]=\mathrm{m}^{-2}$，若 $[\varphi]$ 表示场的 SI 量纲，则 $[\mathcal C_s]=\mathrm{J}\,\mathrm{s}\,\mathrm{m}^{2-D}[\varphi]^{-2}$，保证作用量以焦耳秒计。以下规范代数与 $\mathcal C_s$ 的常数选择无关。
洛伦兹作用量的整体号由所选动力学惯例确定；在这里它不影响规范恒等式及自由场方程。

\begin{proposition}[双迹受限的自由作用量]
对 $T_u^2\Phi=0$ 和 $T_u\varepsilon=0$，上述 Einstein 型作用量在 $\delta\Phi=U\varepsilon$ 下规范不变；紧支变分使边界项消失。
\end{proposition}
\begin{proof}
为核查规范不变性，不必假设 $\mathcal G$ 的普通散度严格为零。利用 $[D_u,u^2]=2U$，先得
\begin{align*}
 D_u\mathcal G
 &=D_u\mathcal F-\frac14D_u(u^2T_u\mathcal F)\\
 &=\left(D_u-\frac12UT_u\right)\mathcal F
   -\frac14u^2D_uT_u\mathcal F\\
 &=-\frac14u^2D_uT_u\mathcal F
 \qquad(T_u^2\Phi=0).
\end{align*}
于是 $D_u\mathcal G$ 的每一项都含显式 $u^2$，而规范参数无迹。在生成多项式的张量缩并配对中，乘 $u^2$ 的伴随操作是取迹 $T_u$，故
\[
 \langle\mathcal G[\Phi], U\varepsilon\rangle_s
 =-\langle D_u\mathcal G[\Phi],\varepsilon\rangle_{s-1}=0.
\]
在双迹受限的变分空间内，$\mathcal G$ 的对称性还允许把 $I_s$ 的一阶变分写成 $\delta I_s=\mathcal C_s\langle\mathcal G[\Phi],\delta\Phi\rangle_s$，从而导出运动方程。\emph{没有双迹约束时，上述消去散度的论证不成立}。这正是原始自由无质量高自旋作用量在 $s\ge4$ 时的结构要求，是由规范不变性决定的约束 \cite{r7-fronsdal-original,r7-covariant-spins}。

\end{proof}

规范场的张量分量个数大于可传播偏振个数。令无质量动量 $p^\mu$ 满足 $p^2=0$ 且 $p^0>0$。固定该动量后，真正作用于偏振张量的紧致旋转群是横向空间的 $SO(D-2)$，不是整个局部 Lorentz 群。采用适当的类 de Donder 规范并商去满足相同波方程的剩余规范变换，可将偏振张量化为横向、对称、无迹秩 $s$ 张量；这个规范固定及约束传播需要用前述 Bianchi 恒等式保证相容。

\begin{proposition}[无质量整数自旋的横向偏振数]\label{prop:ghs-r7-fronsdal-polarizations}
对 $D\ge4$、$s\ge1$ 的上述自由无质量场，局域传播偏振数为
\begin{equation}\label{eq:ghs-r7-fronsdal-physical-count}
 N_{D,s}=\binom{D+s-3}{s}-\binom{D+s-5}{s-2},
\end{equation}
其中当 $s<2$ 时第二项视为零。特别地，$D=4$ 时 $N_{4,s}=2$，对应两个螺旋度 $\pm s$。
\end{proposition}
\begin{proof}
记横向空间维数 $q=D-2\ge2$。对称秩 $s$ 张量的分量数是次数为 $s$ 的 $q$ 变量齐次多项式数目
\[
 \dim\operatorname{Sym}^s\RR^q=\binom{q+s-1}{s}.
\]
迹映射与辅助变量 拉普拉斯算子对应 $T_u: \operatorname{Sym}^s\to\operatorname{Sym}^{s-2}$。在欧氏横向空间，Fischer 分解
\[
 \operatorname{Sym}^s\RR^q
 =\operatorname{Sym}^s_0\RR^q\oplus
 u^2\operatorname{Sym}^{s-2}\RR^q
\]
给出迹映射的满射性及无迹核的维数。因此
\[
 \dim\operatorname{Sym}^s_0\RR^q
 =\binom{q+s-1}{s}-\binom{q+s-3}{s-2}.
\]
代入 $q=D-2$ 即得。对 $q=2$，所有次数 $s\ge1$ 的调和多项式由 $\operatorname{Re}(u_1+\ii u_2)^s$ 与 $\operatorname{Im}(u_1+\ii u_2)^s$ 张成，恰为二维实空间，两个基经过复化正是两个螺旋度偏振。
\end{proof}

这个计数只对具备对应规范约束的\emph{自由无质量}动力学成立。它不等于任意黎曼流形上 $\operatorname{Sym}^s_0(T^*M)$ 的纤维维数，也不能用于一个有质量、没有规范冗余的张量场。

若在一般弯曲背景上试图将 $\partial$ 全部替换为 $\nabla$，情况立即改变。辅助变量需定义在局部正交标架中，协变导数作用于 $u$ 时也包含联络；此时
\[
 [\nabla_\mu,\nabla_\nu]\Phi(x,u)
 =-R^\rho{}_{\sigma\mu\nu}
 u^\sigma\frac{\partial\Phi}{\partial u^\rho}
\]
（指标上下位置随辅助变量采用切向或余切编码而相应转换）。因此原先 $[D_u,U]=\square$、$[T_u,U]=2D_u$ 等平直交换代数不再整体保持：它们的高阶组合会产生曲率乘导数及曲率导数项。一般必须重新求解低阶曲率修正、检查规范约束是否闭合；在特殊的常曲率空间确实存在经过曲率修正的自由高自旋作用量，但这不能证明在任意度规上仍存在同样的无质量规范场。

这一边界给出了这些几何模型两个不同的数学问题：\emph{任意流形上能否定义高自旋张量算子}，以及\emph{该背景是否允许具有正确传播自由度的无质量高自旋规范理论}。前者主要依赖纤维丛和联络，后者还受规范代数、约束传播、主符号与曲率相容性限制。只有明确区分二者，后续的高自旋守恒流和相互作用反常维数才能获得可靠的物理解释。

